% GENERATED FILE. Edit canonical lessons, then run npm run generate:books.
% canonical-source-hash: fnv1a64:d2c7871c0fcbe2ed
% canonical-lessons: TA-C114-natuvil, TA-C114-teru, TA-C114-palam, TA-C114-punka, TA-C114-koyil

\chapter{In the middle, Street, Bridge, Park, Temple}
\label{ch:ta-in-the-middle-street-bridge-park-temple}
\hlchaptermodality{114}

\begin{chapteropening}
\textbf{By the end of this chapter:} \emph{I can say in the middle, a street, a bridge, a park and a temple.}

Complete the last lesson of chapter 114: i can say in the middle, a street, a bridge, a park and a temple.
\end{chapteropening}

\section[\texorpdfstring{\ta{நடுவில்} (naṭuvil)}{நடுவில் (naṭuvil)}]{\ta{நடுவில்} (naṭuvil) — in the middle}

\label{lesson:TA-C114-natuvil}



\begin{quote}
Before the new one: say the Tamil for behind, then the Tamil for in front.
\end{quote}

\subsection*{You'll want to know: \ta{நடுவில்}}
\textbf{\ta{நடுவில்}} — \emph{naṭuvil} — \textquotedblleft{}in the middle\textquotedblright{}.

\textbf{In the family, and next door}

\noindent
\begin{tabularx}{\linewidth}{@{}>{\raggedright\arraybackslash}X>{\raggedright\arraybackslash}X>{\raggedright\arraybackslash}X@{}}
\toprule
\textbf{} & \textbf{word} & \textbf{same root} \\
\midrule
Kannada & \ka{ನಡುವೆ} (\emph{naḍuve}) & yes \\
Telugu & \te{మధ్యలో} (\emph{madhyalō}) &  \\
English & in the middle &  \\
\bottomrule
\end{tabularx}

\begin{grammarlens}[title={What you've built}]
One more word for this chapter.
\end{grammarlens}

\subsection*{Your turn}
\begin{itemize}
\raggedright
  \item \emph{Say it:} \emph{naṭuvil}
  \item \emph{Say it:} \emph{naṭuvil}, once more
  \item \emph{Say it:} say \emph{naṭuvil}
  \item \emph{Recall:} say the Tamil for behind, then the Tamil for in front, then say \emph{naṭuvil} again
\end{itemize}

\begin{tcolorbox}[breakable,colback=teal!4,colframe=teal!35!black,title={Before you move on}]
What does \ta{நடுவில்} mean? (\textquotedblleft{}In the middle\textquotedblright{}.) Say it once more.
\end{tcolorbox}
\section[\texorpdfstring{\ta{தெரு} (teru)}{தெரு (teru)}]{\ta{தெரு} (teru) — a street}

\label{lesson:TA-C114-teru}



\begin{quote}
Before the new one: say the Tamil for in front, then the Tamil for in the middle.
\end{quote}

\subsection*{You'll want to know: \ta{தெரு}}
\textbf{\ta{தெரு}} — \emph{teru} — \textquotedblleft{}a street\textquotedblright{}.

\textbf{In the family, and next door}

\noindent
\begin{tabularx}{\linewidth}{@{}>{\raggedright\arraybackslash}X>{\raggedright\arraybackslash}X>{\raggedright\arraybackslash}X@{}}
\toprule
\textbf{} & \textbf{word} & \textbf{same root} \\
\midrule
Kannada & \ka{ಬೀದಿ} (\emph{bīdi}) &  \\
Telugu & \te{వీధి} (\emph{vīdhi}) &  \\
Malayalam & \ml{തെരുവ്} (\emph{teruvŭ}) & yes \\
English & a street &  \\
\bottomrule
\end{tabularx}

\begin{grammarlens}[title={What you've built}]
One more word for this chapter.
\end{grammarlens}

\subsection*{Your turn}
\begin{itemize}
\raggedright
  \item \emph{Say it:} \emph{teru}
  \item \emph{Say it:} \emph{teru}, once more
  \item \emph{Say it:} say \emph{teru}
  \item \emph{Recall:} say the Tamil for in front, then the Tamil for in the middle, then say \emph{teru} again
\end{itemize}

\begin{tcolorbox}[breakable,colback=teal!4,colframe=teal!35!black,title={Before you move on}]
What does \ta{தெரு} mean? (\textquotedblleft{}A street\textquotedblright{}.) Say it once more.
\end{tcolorbox}
\section[\texorpdfstring{\ta{பாலம்} (pālam)}{பாலம் (pālam)}]{\ta{பாலம்} (pālam) — a bridge}

\label{lesson:TA-C114-palam}



\begin{quote}
Before the new one: say the Tamil for in the middle, then the Tamil for a street.
\end{quote}

\subsection*{You'll want to know: \ta{பாலம்}}
\textbf{\ta{பாலம்}} — \emph{pālam} — \textquotedblleft{}a bridge\textquotedblright{}.

\textbf{In the family, and next door}

\noindent
\begin{tabularx}{\linewidth}{@{}>{\raggedright\arraybackslash}X>{\raggedright\arraybackslash}X>{\raggedright\arraybackslash}X@{}}
\toprule
\textbf{} & \textbf{word} & \textbf{same root} \\
\midrule
Kannada & \ka{ಸೇತುವೆ} (\emph{sētuve}) &  \\
Telugu & \te{వంతెన} (\emph{vantena}) &  \\
Malayalam & \ml{പാലം} (\emph{pālaṁ}) & yes \\
Hindi (neighbour) & \dv{पुल} (\emph{pul}) &  \\
English & a bridge &  \\
\bottomrule
\end{tabularx}

\begin{grammarlens}[title={What you've built}]
One more word for this chapter.
\end{grammarlens}

\subsection*{Your turn}
\begin{itemize}
\raggedright
  \item \emph{Say it:} \emph{pālam}
  \item \emph{Say it:} \emph{pālam}, once more
  \item \emph{Say it:} say \emph{pālam}
  \item \emph{Recall:} say the Tamil for in the middle, then the Tamil for a street, then say \emph{pālam} again
\end{itemize}

\begin{tcolorbox}[breakable,colback=teal!4,colframe=teal!35!black,title={Before you move on}]
What does \ta{பாலம்} mean? (\textquotedblleft{}A bridge\textquotedblright{}.) Say it once more.
\end{tcolorbox}
\section[\texorpdfstring{\ta{பூங்கா} (pūṅkā)}{பூங்கா (pūṅkā)}]{\ta{பூங்கா} (pūṅkā) — a park}

\label{lesson:TA-C114-punka}



\begin{quote}
Before the new one: say the Tamil for a street, then the Tamil for a bridge.
\end{quote}

\subsection*{You'll want to know: \ta{பூங்கா}}
\textbf{\ta{பூங்கா}} — \emph{pūṅkā} — \textquotedblleft{}a park\textquotedblright{}.

\textbf{In the family, and next door}

\noindent
\begin{tabularx}{\linewidth}{@{}>{\raggedright\arraybackslash}X>{\raggedright\arraybackslash}X@{}}
\toprule
\textbf{} & \textbf{word} \\
\midrule
Kannada & \ka{ಉದ್ಯಾನ} (\emph{udyāna}) \\
Telugu & \te{పార్కు} (\emph{pārku}) \\
English & a park \\
\bottomrule
\end{tabularx}

\begin{grammarlens}[title={What you've built}]
One more word for this chapter.
\end{grammarlens}

\subsection*{Your turn}
\begin{itemize}
\raggedright
  \item \emph{Say it:} \emph{pūṅkā}
  \item \emph{Say it:} \emph{pūṅkā}, once more
  \item \emph{Say it:} say \emph{pūṅkā}
  \item \emph{Recall:} say the Tamil for a street, then the Tamil for a bridge, then say \emph{pūṅkā} again
\end{itemize}

\begin{tcolorbox}[breakable,colback=teal!4,colframe=teal!35!black,title={Before you move on}]
What does \ta{பூங்கா} mean? (\textquotedblleft{}A park\textquotedblright{}.) Say it once more.
\end{tcolorbox}
\section[\texorpdfstring{\ta{கோயில்} (kōyil)}{கோயில் (kōyil)}]{\ta{கோயில்} (kōyil) — a temple}

\label{lesson:TA-C114-koyil}



\begin{quote}
Before the new one: say the Tamil for a bridge, then the Tamil for a park.

Name a Tamil month, \textbf{\ta{சித்திரை}} or \textbf{\ta{தை}}, then give a \textbf{\ta{தேதி}}, with the date ending \emph{-ām}.
\end{quote}

\subsection*{You'll want to know: \ta{கோயில்}}
\textbf{\ta{கோயில்}} — \emph{kōyil} — \textquotedblleft{}a temple\textquotedblright{}.

\textbf{In the family, and next door}

\noindent
\begin{tabularx}{\linewidth}{@{}>{\raggedright\arraybackslash}X>{\raggedright\arraybackslash}X@{}}
\toprule
\textbf{} & \textbf{word} \\
\midrule
Kannada & \ka{ದೇವಸ್ಥಾನ} (\emph{dēvasthāna}) \\
Telugu & \te{గుడి} (\emph{guḍi}) \\
Malayalam & \ml{ക്ഷേത്രം} (\emph{kṣētraṁ}) \\
Hindi (neighbour) & \dv{मंदिर} (\emph{mandir}) \\
English & a temple \\
\bottomrule
\end{tabularx}

\begin{grammarlens}[title={What you've built}]
One more word for this chapter.
\end{grammarlens}

\subsection*{Your turn}
\begin{itemize}
\raggedright
  \item \emph{Say it:} \emph{kōyil}
  \item \emph{Say it:} \emph{kōyil}, once more
  \item \emph{Say it:} say \emph{kōyil}
  \item \emph{Recall:} say the Tamil for a bridge, then the Tamil for a park, then say \emph{kōyil} again
\end{itemize}

\begin{tcolorbox}[breakable,colback=teal!4,colframe=teal!35!black,title={Before you move on}]
What does \ta{கோயில்} mean? (\textquotedblleft{}A temple\textquotedblright{}.) Say it once more.
\end{tcolorbox}
